\subsection{HIGHEST WEIGHT OF AN IRREDUCIBLE REPRESENTATION}

Associate to any representation \(\tau : \mathrm{G} \to \mathrm{GL}(\mathrm{V})\) the homomorphism \(\mathrm{L}(\tau)_{(\mathbf{C})}\) from the \(\mathbf{C}\) -Lie algebra \(\mathfrak{g}_{\mathbf{C}}\) to \(\operatorname{End}(\mathrm{V})\) extending the linear representation \(\mathrm{L}(\tau)\) of \(\mathfrak{g}\) on the real vector space underlying V (Chap. III, §3, no. 11). By Prop. 7 of §4, no. 3, the map \(\delta\) from \(\mathrm{X}(\mathrm{T})\) to \(\operatorname{Hom}_{\mathbf{C}}(\mathfrak{t}_{\mathbf{C}}, \mathbf{C}) = \mathfrak{t}_{\mathbf{C}}^{*}\) induces a bijection from the set of weights of \(\tau\) relative to T to the set of weights of \(\mathrm{L}(\tau)_{(\mathbf{C})}\) relative to the Cartan subalgebra \(\mathfrak{t}_{\mathbf{C}}\) of \(\mathfrak{g}_{\mathbf{C}}\) .

Lemma 2. Let \(\varphi\) be a linear representation of the complex Lie algebra \(\mathfrak{g}_{\mathbf{C}}\) on a finite dimensional complex vector space V. There exists a representation \(\tau\) of G on V such that \(\mathrm{L}(\tau)_{(\mathbf{C})} = \varphi\) if and only if \(\varphi\) is semi-simple and the weights of \(\mathbf{t}_{\mathbf{C}}\) on V belong to \(\delta (\mathrm{X(T)})\) .

If there exists a representation \(\tau\) of G such that \(\mathrm{L}(\tau)_{(\mathbf{C})} = \varphi\) , then \(\varphi\) is semi-simple because G is connected and \(\tau\) is semi-simple (Chap. III, §6, no. 5, Cor. 2 of Prop. 13), and the weights of \(\mathfrak{t}_{\mathbf{C}}\) on V belong to the image of \(\delta\) . Thus, the condition is necessary; we shall prove that it is sufficient. If \(\varphi\) is semi-simple, V is the direct sum of the \(V_{\mu}(t_{\mathbf{C}})\) , where \(\mu\) belongs to the set of weights of \(\mathfrak{t}_{\mathbf{C}}\) on V (Chap. VII, §2, no. 4, Cor. 3 of Th. 2); if all the weights belong to the image of \(\delta\) , there exists a representation \(\tau_{\mathrm{T}}\) of T on V such that \(\mathrm{L}(\tau_{\mathrm{T}})_{(\mathbf{C})} = \varphi |t_{\mathbf{C}}\) : indeed, it suffices to put \(\tau_{\mathrm{T}}(t)v = t^{\lambda}v\) for \(t \in \mathrm{T}\) and \(v \in V_{\delta(\lambda)}(t_{\mathbf{C}})\) . The lemma now follows from Prop. 8 of §2, no. 6.

THEOREM 1. a) Let \(\tau : \mathbf{G} \to \mathbf{GL}(\mathbf{V})\) be an irreducible representation of \(\mathbf{G}\) . Then the set of weights of \(\tau\) (relative to \(\mathbf{T}\) ) has a largest element \(\lambda\) , which is dominant, and the space \(\mathrm{V}_{\lambda}(\mathrm{T})\) is of dimension 1.

\begin{enumerate}
\def\labelenumi{\alph{enumi})}
\setcounter{enumi}{1}
\item
  Two irreducible representations of G are equivalent if and only if their highest weights are equal.
\item
  For every dominant element \(\lambda\) of X(T), there exists an irreducible representation of G of highest weight \(\lambda\) .
\end{enumerate}

By Lemma 2, the equivalence classes of irreducible representations of G correspond bijectively to the classes of finite dimensional irreducible representations of g whose weights belong to \(\delta(\mathrm{X}(\mathrm{T}))\) .

Denote by \(\mathcal{G}\mathfrak{g}_{\mathbf{C}}\) the centre and by \(\mathcal{D}\mathfrak{g}_{\mathbf{C}}\) the derived Lie algebra of \(\mathfrak{g}_{\mathbf{C}}\) , so that \(\mathfrak{g}_{\mathbf{C}} = \mathcal{G}\mathfrak{g}_{\mathbf{C}} \oplus \mathcal{D}\mathfrak{g}_{\mathbf{C}}\) . For every bilinear form \(\mu\) on \(\mathfrak{t}_{\mathbf{C}} \cap \mathcal{D}\mathfrak{g}_{\mathbf{C}}\) , denote by \(\mathrm{E}(\mu)\) the simple \(\mathcal{D}\mathfrak{g}_{\mathbf{C}}\) -module introduced in Chap. VIII, §6, no. 3; for every linear form \(\nu\) on \(\mathcal{G}\mathfrak{g}_{\mathbf{C}}\) , denote by \(\mathbf{C}(\nu)\) the 1-dimensional \(\mathcal{G}\mathfrak{g}_{\mathbf{C}}\) -module over \(\mathbf{C}\) associated to it. Then the \(\mathfrak{g}_{\mathbf{C}}\) -modules \(\mathbf{C}(\nu) \otimes \mathrm{E}(\mu)\) are simple, and by Chap. VIII, §7, no. 2, Cor. 2 of Th. 1 and Algebra, Chap. VIII, §11, no. 1, Th. 1, every finite dimensional simple \(\mathfrak{g}_{\mathbf{C}}\) -module is isomorphic to one of the modules \(\mathbf{C}(\nu) \otimes \mathrm{E}(\mu)\) ; moreover (loc. cit.) \(\mathbf{C}(\nu) \otimes \mathrm{E}(\mu)\) is finite dimensional if and only if \(\mu(H_{\alpha})\) is a positive integer for every simple root \(\alpha\) . If we denote by \(\nu + \mu\) the linear form on \(\mathfrak{t}_{\mathbf{C}}\) that induces \(\nu\) on \(\mathcal{G}\mathfrak{g}_{\mathbf{C}}\) and \(\mu\) on \(\mathfrak{t}_{\mathbf{C}} \cap \mathcal{D}\mathfrak{g}_{\mathbf{C}}\) , then \((\nu + \mu)(H_{\alpha}) = \mu(H_{\alpha})\) ; moreover, the weights of \(\mathbf{C}(\nu) \otimes \mathrm{E}(\mu)\) are the \(\nu + \lambda\) , where \(\lambda\) is any weight of \(\mathrm{E}(\mu)\) , and are thus of the form \(\nu + \mu - \theta\) , with \(\theta \in \delta(X_{+})\) (Chap. VIII, §6, no. 2, Lemma 2).

We conclude that the g-module \(\mathbf{C}(\nu)\otimes\mathrm{E}(\mu)\) is finite dimensional if and only if \((\nu+\mu)(H_{\alpha})\) is a positive integer for every simple root \(\alpha\) , and that its weights belong to \(\delta(\mathrm{X}(\mathrm{T}))\) if and only if \(\nu+\mu\) belongs to \(\delta(\mathrm{X}(\mathrm{T}))\) . The conjunction of these two conditions means that \(\nu+\mu\) belongs to \(\delta(\mathrm{X}_{++})\) ; in that case, \(\nu+\mu\) is the highest weight of \(\mathbf{C}(\nu)\otimes\mathrm{E}(\mu)\) . Thus, we have constructed for every dominant weight \(\lambda\) of X(T) an irreducible representation of G of highest weight \(\lambda\) , and have thus obtained, up to equivalence, all the irreducible representations of G. Since the vectors of weight \(\nu+\mu\) in \(\mathbf{C}(\nu)\otimes\mathrm{E}(\mu)\) form a subspace of dimension 1, the proof is complete.

COROLLARY. The group G has a (finite dimensional) faithful linear representation.

Observe first that every element of X(T) is equal to the difference of two dominant weights: more precisely, let \(\varpi\) be an element of \(X_{++}\) such that \(\langle\varpi,K_{\alpha}\rangle>0\) for every simple root \(\alpha\) ; for all \(\lambda\in X(T)\) there exists a positive integer n such that \(\langle\lambda+n\varpi,K_{\alpha}\rangle\geq0\) for every simple root \(\alpha\) , that is (no. 1, Lemma 1) \(\lambda+n\varpi\in X_{++}\) .

Consequently, there exists a finite family \((\lambda_{i})_{i\in I}\) of elements of \(X_{++}\) generating the Z-module \(\mathrm{X(T)}\) . For \(i\in I\) , let \(\tau_{i}\) be an irreducible representation of G of highest weight \(\lambda_{i}\) (Th. 1); let the representation \(\tau\) be the direct sum of the \(\tau_{i}\) . By construction the set \(\mathrm{P}(\tau,\mathrm{T})\) of weights of \(\tau\) (relative to T) generates the Z-module \(\mathrm{X(T)}\) . It now follows from Prop. 6 of §4, no. 3 that the homomorphism \(\tau\) is injective, hence the corollary.

Remarks. 1) Let \(\mathfrak{n}_{+}\) be the subalgebra of \(\mathfrak{g}_{\mathbf{C}}\) that is the sum of the \(\mathfrak{g}^{\alpha}\) for \(\alpha > 0\) . Let \(\tau: \mathrm{G} \to \mathrm{GL}(\mathrm{V})\) be an irreducible representation, \(\lambda \in \mathrm{X}_{++}\) its highest weight and \(\tau': \mathfrak{g}_{\mathbf{C}} \to \mathfrak{gl}(\mathrm{V})\) the representation induced by \(\tau\) . Then \(\mathrm{V}_{\lambda}(\mathrm{T})\) is the subspace consisting of the vectors \(v\) in \(\mathrm{V}\) such that \(\tau'(x)v = 0\) for all \(x \in \mathfrak{n}_{+}\) .

Indeed, this follows from the corresponding statement for the \(\mathfrak{g}_{\mathbf{C}}\) -modules \(\mathbf{C}(\nu) \otimes \mathrm{E}(\mu)\) (Chap. VIII, §6, no. 2, Prop. 3).

\begin{enumerate}
\def\labelenumi{\arabic{enumi})}
\setcounter{enumi}{1}
\item
  Let \(\Theta(\mathrm{G})\) be the algebra of continuous representative functions on \(\mathrm{G}\) with values in \(\mathbf{C}\) (Algebra, Chap. VIII). Let \(\mathrm{G}\) operate on \(\Theta(\mathrm{G})\) by left and right translations. For each \(\lambda \in \mathrm{X}_{++}\) , let \((\mathrm{V}_{\lambda}, \tau_{\lambda})\) be an irreducible representation of \(\mathrm{G}\) of highest weight \(\lambda\) (Th. 1), and \((\mathrm{V}_{\lambda}^{*}, \check{\tau}_{\lambda})\) the contragredient representation (Chap. III, §3, no. 11); by Spectral Theories, the representation of \(\mathrm{G} \times \mathrm{G}\) on \(\Theta(\mathrm{G})\) is isomorphic to the direct sum of the representations \((\mathrm{V}_{\lambda} \otimes \mathrm{V}_{\lambda}^{*}, \tau_{\lambda} \otimes \check{\tau}_{\lambda})\) for all \(\lambda\) in \(\mathrm{X}_{++}\) . Remark 1 now implies the following statement: Let \(\lambda \in \mathrm{X}_{++}\) , and let \(\mathrm{E}_{\lambda}\) be the vector subspace of \(\Theta(\mathrm{G})\) consisting of the continuous representative functions \(f\) on \(\mathrm{G}\) such that \(f(gt) = \lambda(t)^{-1} f(g)\) for all \(g \in \mathrm{G}\) and all \(t \in \mathrm{T}\) , and such that \(f * x = 0\) for all \(x \in \mathfrak{n}_{-} = \bigoplus_{\alpha < 0} \mathfrak{g}^{\alpha}\) . Then \(\mathrm{E}_{\lambda}\) is stable under left translations, and the representation of \(\mathrm{G}\) on \(\mathrm{E}_{\lambda}\) by left translations is irreducible of highest weight \(\lambda\) .
\item
  Let \(\tau : \mathrm{G} \to \mathbf{GL}(\mathrm{V})\) be an irreducible representation. There exists an element \(\nu\) of \(\mathrm{X}(\mathrm{C}(\mathrm{G}))\) such that \(\tau(s)v = \nu(s)v\) for all \(s \in \mathrm{C}(\mathrm{G}), v \in \mathrm{V}\) : indeed, \(\tau(\mathrm{C}(\mathrm{G}))\) is contained in the commutant of \(\tau(\mathrm{G})\) , which is equal to \(\mathbf{C}^*.1_{\mathrm{V}}\) (Algebra, Chap. VIII, §3, no. 2, Th. 1). For every weight \(\lambda\) of \(\tau\) , the restriction of \(\lambda\) to \(\mathrm{C}(\mathrm{G})\) is equal to \(\nu\) .
\item
  The definitions and statements of Chap. VIII, §7, nos. 2 to 5 can be generalized without difficulty to the present situation; we leave the details to the reader.
\end{enumerate}

PROPOSITION 1. Let \(\tau : G \to \mathbf{GL}(V)\) be an irreducible representation of \(G\) of highest weight \(\lambda \in X_{++}\) . Let \(m\) be the integer \(\sum_{\alpha \in R_+} \langle \lambda, K_\alpha \rangle\) , and let \(w_0\) be the element of the Weyl group such that \(w_0(R_+) = R_-\) (Chap. VI, § 6, Cor. 3 of Prop. 17). There are three possible cases:

\begin{enumerate}
\def\labelenumi{\alph{enumi})}
\item
  \(w_{0}(\lambda) = -\lambda\) and m is even. Then there exists a G-invariant separating symmetric bilinear form on V; the representation \(\tau\) is of real type (Appendix II).
\item
  \(w_{0}(\lambda) \neq -\lambda\) . Every G-invariant bilinear form on V is zero; the representation \(\tau\) is of complex type (loc. cit.).
\item
  \(w_{0}(\lambda) = -\lambda\) and m is odd. Then there exists a G-invariant separating alternating bilinear form on V; the representation \(\tau\) is of quaternionic type (loc. cit.).
\end{enumerate}

If the restriction of \(\tau\) to \(\mathrm{C}(\mathrm{G})_0\) is not trivial, we are in case \(b\) ).

A bilinear form B on V is invariant under G if and only if it is invariant under \(g_{C}\) (Chap. III, §6, no. 5, Cor. 3). Thus, if G is semi-simple, the proposition follows from Chap. VIII, §7, no. 5, Prop. 12 and Prop. 3 of Appendix II.

In the general case, put \(\mathrm{C}(\mathrm{G})_{0} = \mathrm{S}\) , and identify \(\mathrm{X}(\mathrm{T}/\mathrm{S})\) with a subgroup of \(\mathrm{X}(\mathrm{T})\) (stable under W). If \(\tau(\mathrm{S}) = \{1_{\mathrm{V}}\}\) , \(\tau\) induces by passage to the quotient a representation \(\tau' : G/S \to GL(V)\) of highest weight \(\lambda\) ; in this case the proposition follows from the preceding, applied to G/S.

Assume that \(\tau (\mathrm{S})\neq \{1_{\mathrm{V}}\}\) . There exists a non-zero element \(\nu\) of \(\mathrm{X(S)}\) such that \(\tau (s) = \nu (s)_{\mathrm{V}}\) for all \(s\in \mathrm{S}\) (Remark 3). Then \(\nu\) is the image of \(\lambda\) under the restriction homomorphism \(\mathrm{X}(\mathrm{T})\to\mathrm{X}(\mathrm{S})\) ; since W operates trivially on \(\mathrm{X}(\mathrm{S})\) , the equality \(w_{0}(\lambda)=-\lambda\) implies that \(\nu=-\nu\) , which is impossible: hence \(w_{0}(\lambda)\neq-\lambda\) . On the other hand, if B is a G-invariant bilinear form on V, we have, for all x,y in V and s in S,

\[
\mathrm{B} (\nu (s) x, \nu (s) y) = \mathrm{B} (x, y) = \nu (s) ^ {2} \mathrm{B} (x, y)
\]

which implies that B = 0, hence the proposition.

Let \(\mathfrak{S}_{\mathbf{R}}(\mathrm{G})\) be the set of classes of irreducible continuous representations of G on finite dimensional real vector spaces. Prop. 1 and the results of Appendix II give a bijection \(\Phi:X_{++}/\Sigma\to\mathfrak{S}_{\mathbf{R}}(\mathrm{G})\) , where \(\Sigma\) denotes the subgroup \(\{1,-w_{0}\}\) of \(\operatorname{Aut}(\mathrm{X}(\mathrm{T}))\) . More precisely, let \(\lambda\in X_{++}\) , and let \(E_{\lambda}\) be a representation of G of highest weight \(\lambda\) ; then

\[
\begin{array}{c} \Phi (\{\lambda , - w _ {0} (\lambda) \}) = \mathrm{E} _ {\lambda [ \mathbf {R} ]} \text {if} \lambda \neq - w _ {0} (\lambda) \text {or if} \sum_ {\alpha \in \mathrm{R} _ {+}} \langle \lambda , K _ {\alpha} \rangle \notin 2 \mathbf {Z} \\ \Phi (\{\lambda , - w _ {0} (\lambda) \}) = \mathrm{E} _ {\lambda} ^ {\prime} \text {if} \lambda = - w _ {0} (\lambda) \text {and} \sum_ {\alpha \in \mathrm{R} _ {+}} \langle \lambda , K _ {\alpha} \rangle \in 2 \mathbf {Z} \end{array}
\]

where \(E_{\lambda}^{\prime}\) is an R-structure on \(E_{\lambda}\) invariant under G.

